\chapter{Zawieszenie, twierdzenie Freudenthala i homotopia stabilna}
\label{ch:stabilna-v8}

\section{Motywacja: co pozostaje po wielokrotnym zawieszaniu?}
Dla każdej liczby $n\ge1$ wiemy, że
$\pi_n(S^n)\cong\Z$. Dla grup
$\pi_{n+k}(S^n)$ przy ustalonym $k\ge1$ wyniki
są trudniejsze, lecz po dostatecznym zwiększeniu $n$
grupy przestają się zmieniać. Ta stabilizacja jest
konsekwencją twierdzenia Freudenthala.

W tym rozdziale konsekwentnie używamy \emph{zredukowanego}
zawieszenia
$\Sigma X=S^1\wedge X$ dla CW-kompleksu z punktem
bazowym. Dla $X=S^n$ jest ono homeomorficzne z $S^{n+1}$.
Przyjmujemy standardowe punkty bazowe włączone jako
podkompleksy, aby ilorazy i pary były kofibracjami.

\begin{definicja}[Zawieszenie odwzorowania]
Dla punktowanego odwzorowania $f:X\to Y$
definiujemy $\Sigma f=\id_{S^1}\wedge f:
\Sigma X\to\Sigma Y$.
Jeżeli $f:(S^q,*)\to(X,x_0)$, to
$\Sigma f:S^{q+1}\to\Sigma X$ wyznacza
\emph{homomorfizm zawieszenia}
$E:\pi_q(X)\to\pi_{q+1}(\Sigma X)$.
Dla $q=1$ odwzorowanie to rozumiemy jako
homomorfizm z grupy,
niekoniecznie abelowej, do abelowej grupy $\pi_2$.
\end{definicja}

\begin{stwierdzenie}[Naturalność]
Jeżeli $f:X\to Y$ jest punktowane, to dla $q\ge1$
zachodzi $(\Sigma f)_*\circ E_X=E_Y\circ f_*$.
\end{stwierdzenie}
\begin{proof}
Dla reprezentanta $\alpha:S^q\to X$ oba złożenia są
reprezentowane przez to samo odwzorowanie
$\Sigma(f\circ\alpha)=(\Sigma f)\circ(\Sigma\alpha)$.
\end{proof}

\section{Spójność przestrzeni i ekscyzja homotopijna}
\begin{definicja}[$r$-spójność]
Niech $r\ge0$. Punktowaną przestrzeń $X$ nazywamy
\emph{$r$-spójną}, jeśli jest łukowo spójna i
$\pi_i(X)=0$ dla $1\le i\le r$.
Parę punktowaną $(Y,A)$ nazywamy $r$-spójną,
jeśli $\pi_i(Y,A)=0$ dla $1\le i\le r$
oraz $\pi_0(A)\to\pi_0(Y)$ jest surjekcją.
Dla $r=0$ oznacza to tylko warunek na składowe.
\end{definicja}

\begin{lemat}[Stożek i względne grupy homotopii]
\label{lem:cone-relative-v8}
Niech $X$ będzie $r$-spójnym CW-kompleksem,
$r\ge0$. Wtedy para $(CX,X)$ jest $(r+1)$-spójna,
a dla $q\ge2$ odwzorowanie brzegowe
$\partial:\pi_{q+1}(CX,X)\to\pi_q(X)$
jest izomorfizmem.
\end{lemat}
\begin{proof}
Stożek jest ściągalny, więc $\pi_q(CX)=0$ dla
$q\ge1$. Fragment długiego ciągu homotopii pary ma postać
$\pi_{q+1}(CX)\to\pi_{q+1}(CX,X)
\xrightarrow{\partial}\pi_q(X)\to\pi_q(CX)$.
Skrajne grupy są zerowe, stąd izomorfizm.
Dla $q\le r$ zeruje się też $\pi_{q+1}(CX,X)$;
warunki w niskich stopniach wynikają z łączności $X$.
\end{proof}

\begin{twierdzenie}[Ekscyzja homotopijna Blakersa--Masseya]
\label{thm:bm-v8}
Niech $Z=A\cup B$ będzie triadą CW, gdzie
$C=A\cap B$ jest niepuste i łukowo spójne.
Jeżeli pary $(A,C)$ i $(B,C)$ są odpowiednio
$m$-spójna i $n$-spójna, $m,n\ge1$, to homomorfizm
$\pi_q(A,C)\to\pi_q(Z,B)$ jest izomorfizmem dla
$q<m+n$ i surjekcją dla $q=m+n$.
\end{twierdzenie}
\begin{uwaga}[Status dowodu ekscyzji homotopijnej]
Jest to głębokie twierdzenie i \emph{nie} wynika ze
zwykłej ekscyzji homologicznej. Pełny dowód można
przeprowadzić przez redukcję do triad CW oraz
usuwanie przecięć obrazów kostek z komórkami
(ogólna pozycja). Wykorzystujemy je jako
wcześniej jawnie sformułowany wynik pomocniczy;
dowód Freudenthala poniżej jest pełny \emph{względem}
tego twierdzenia, nie stanowi niezależnego dowodu
Blakersa--Masseya.
\end{uwaga}

\begin{twierdzenie}[Freudenthal]
\label{thm:freudenthal-v8}
Niech $X$ będzie $r$-spójnym, dobrze bazowanym
CW-kompleksem, gdzie $r\ge0$. Wówczas homomorfizm
\[
 E:\pi_q(X)\longrightarrow\pi_{q+1}(\Sigma X)
\]
jest izomorfizmem dla $1\le q<2r+1$ oraz
surjekcją dla $q=2r+1$.
\end{twierdzenie}
\begin{proof}
Przedstawmy zawieszenie jako sumę dwóch
zredukowanych stożków $C_+X$ i $C_-X$,
których przecięcie ma typ homotopijny $X$.
Oba stożki są ściągalne. Pary $(C_+X,X)$
i $(C_-X,X)$ są $(r+1)$-spójne na mocy
Lematu~\ref{lem:cone-relative-v8}.
Ekscyzja homotopijna dla tej triady daje
$\pi_{q+1}(C_+X,X)\to
\pi_{q+1}(\Sigma X,C_-X)$
izomorfizm, gdy $q+1<2(r+1)$, i surjekcję,
gdy $q+1=2(r+1)$.

Z długich ciągów homotopii par i ściągalności
stożków otrzymujemy naturalne identyfikacje
$\pi_{q+1}(C_+X,X)\cong\pi_q(X)$ oraz
$\pi_{q+1}(\Sigma X,C_-X)\cong
\pi_{q+1}(\Sigma X)$ dla $q\ge2$.
Homomorfizm pomiędzy tymi dwiema grupami
przyporządkowuje reprezentantowi $f:S^q\to X$
odwzorowanie $\Sigma f:S^{q+1}\to\Sigma X$,
co widać przez potraktowanie reprezentanta względnego
jako odwzorowania $(D^{q+1},S^q)\to(C_+X,X)$
i sklejenie drugiego stożka.
Zatem wskazana ekscyzja jest dokładnie homomorfizmem $E$.
Daje to podany zakres izomorfizmu i surjekcji.
Pozostał przypadek $q=1$. Jeżeli $r\ge1$, obie
grupy są zerowe: $X$ jest jednospójna, a $\Sigma X$
jest co najmniej $2$-spójna. Istotnie, $\Sigma X$ jest
jednospójna, ma $H_2(\Sigma X;\Z)\cong H_1(X;\Z)=0$,
a więc $\pi_2(\Sigma X)=0$ z twierdzenia Hurewicza. Jeżeli $r=0$, przestrzeń $X$
jest łukowo spójna i $\Sigma X$ jest jednospójna.
Twierdzenie Hurewicza, izomorfizm zawieszenia w homologii
oraz twierdzenie o abelianizacji grupy podstawowej dają
naturalne identyfikacje
\[
\pi_2(\Sigma X)\cong H_2(\Sigma X;\Z)
\cong H_1(X;\Z)\cong\pi_1(X)_{\mathrm{ab}}.
\]
Aby sprawdzić, że to właśnie $E$, weźmy pętlę
$\alpha:S^1\to X$; obraz klasy $E[\alpha]$ przez
homomorfizm Hurewicza jest zawieszeniem cyklu
reprezentującego $[\alpha]$. Wynika to z naturalności
homomorfizmu Hurewicza oraz izomorfizmu zawieszenia dla
$S^1$ i $X$. Zatem $E$ jest odwzorowaniem abelianizacji,
w szczególności surjekcją, czego wymaga teza.
\end{proof}

\begin{center}
\begin{tikzpicture}[>=Stealth,font=\small]
\node[draw,rounded corners,minimum width=30mm,minimum height=11mm] (X) at (0,0) {$X$};
\node[draw,rounded corners,minimum width=30mm,minimum height=11mm] (SX) at (4.3,0) {$\Sigma X$};
\node[draw,rounded corners,minimum width=30mm,minimum height=11mm] (SSX) at (8.6,0) {$\Sigma^2 X$};
\draw[->,thick] (X)-- node[above]{zawieszenie} (SX);
\draw[->,thick] (SX)-- node[above]{zawieszenie} (SSX);
\node[font=\footnotesize] at (4.3,-1.05)
{Po dostatecznie wielu krokach grupy $\pi_{q+k}(\Sigma^kX)$ stabilizują się.};
\end{tikzpicture}
\end{center}

\begin{przyklad}[Ostry brzeg zakresu]
Sfera $S^2$ jest $1$-spójna, więc
$E:\pi_2(S^2)\to\pi_3(S^3)$ jest izomorfizmem,
ale dla $q=3$ twierdzenie gwarantuje już tylko surjekcję
$E:\pi_3(S^2)\to\pi_4(S^3)$.
Pokażemy dalej, że jest to redukcja $\Z\to\Z/2$.
W szczególności zastąpienie w twierdzeniu
warunku $q<2r+1$ przez $q\le2r+1$
byłoby nieprawdziwe.
\end{przyklad}

\begin{przyklad}[Zawieszenie abelianizuje grupę podstawową]
\label{prz:freudenthal-klein-abel-v8}
Rozważmy butelkę Kleina $K$ z
Przykładu~\ref{prz:fib-klein-monodromia}.
Znamy prezentację
\[
 \pi_1(K)=\langle a,b\mid aba^{-1}=b^{-1}\rangle.
\]
W abelianizacji elementy $a,b$ komutują,
więc relacja upraszcza się do $b=b^{-1}$,
czyli $2[b]=0$.
Nie pojawia się relacja na $[a]$:
homomorfizm $a\mapsto1$, $b\mapsto0$
na $\Z$ respektuje relację wyjściową.
Dlatego
\[
 \pi_1(K)_{\mathrm{ab}}
 \cong\Z\oplus\Z/2\Z.
\]
Ponieważ $K$ jest łukowo spójna,
$\Sigma K$ jest jednospójna, a twierdzenie Hurewicza
oraz izomorfizm zawieszenia w homologii dają
\[
 \pi_2(\Sigma K)
 \cong H_2(\Sigma K;\Z)
 \cong H_1(K;\Z)
 \cong\Z\oplus\Z/2\Z.
\]
Przy tych identyfikacjach
$E:\pi_1(K)\to\pi_2(\Sigma K)$
jest dokładnie odwzorowaniem abelianizacji.
Jest \emph{surjektywne}, lecz nie jest różnowartościowe:
jego jądro zawiera komutatory, na przykład
$[a,b]=aba^{-1}b^{-1}=b^{-2}$,
który nie jest neutralny w wyjściowej grupie
(z produktu półprostego $\Z\rtimes_{-1}\Z$
wynika, że $b$ ma nieskończony rząd).
To ilustruje ostrość granicy $q=1$, $r=0$
w twierdzeniu Freudenthala.

Dla porównania na torusie $T^2$ mamy
$\pi_1(T^2)=\Z^2$ już abelową, a więc
$E:\pi_1(T^2)\to\pi_2(\Sigma T^2)\cong\Z^2$
jest izomorfizmem. Warunek
$q=2r+1$ daje gwarancję surjektywności,
nie stwierdza, że w każdym przykładzie
izomorfizm musi zawieść.
\end{przyklad}

\section{Stabilne grupy homotopii sfer}
\begin{definicja}[Stabilna grupa homotopii]
Dla dobrze bazowanego CW-kompleksu $X$ i $k\in\Z$
definiujemy
\[
\pi_k^{\mathrm s}(X)=
\varinjlim_{j\ge\max(1,1-k)}
\pi_{k+j}(\Sigma^jX),
\]
gdzie strzałkami układu są homomorfizmy zawieszenia.
Dolna granica indeksowania zapewnia, że wszystkie
występujące grupy homotopii mają dodatni stopień;
odrzucenie skończenie wielu początkowych wyrazów
układu nie zmienia jego granicy prostej.
W szczególności dla $k\ge0$ kładziemy
$\pi_k^{\mathrm s}=\pi_k^{\mathrm s}(S^0)
=\varinjlim_{n}\pi_{n+k}(S^n)$.
Grupy $\pi_k^{\mathrm s}$ nazywamy \emph{stabilnymi
pniami homotopii sfer} (dla $k\ge0$).
\end{definicja}

\begin{twierdzenie}[Stabilizacja w ustalonym stopniu]
\label{thm:stable-range-v8}
Niech $X$ będzie niepustym, dobrze bazowanym
CW-kompleksem. Dla każdego ustalonego $k\ge0$
odwzorowania
$\pi_{k+j}(\Sigma^j X)\to
\pi_{k+j+1}(\Sigma^{j+1}X)$
są izomorfizmami dla dostatecznie dużych $j$.
Dla sfer $S^n$ i $k\ge0$ wystarcza $n\ge k+2$;
przy tym założeniu zawieszenie
$\pi_{n+k}(S^n)\to\pi_{n+k+1}(S^{n+1})$
jest izomorfizmem.
\end{twierdzenie}
\begin{proof}
Po pierwszym zawieszeniu $\Sigma X$ jest łukowo spójne.
Ponadto każde kolejne zawieszenie przestrzeni
łukowo spójnej podwyższa jej spójność o co najmniej
jeden. Rzeczywiście, zawieszenie przestrzeni
łukowo spójnej jest jednospójne, a izomorfizm
zawieszenia w homologii oraz twierdzenie Hurewicza
pozwalają kolejno wyzerować grupy homotopii
w niższych stopniach. W konsekwencji, dla $j\ge1$,
przestrzeń $\Sigma^jX$ jest co najmniej
$(j-1)$-spójna, także gdy pierwotna $X$
nie była łukowo spójna.
Stosujemy Twierdzenie~\ref{thm:freudenthal-v8}
z $r=j-1$ do grupy o indeksie $q=k+j$.
Wymagamy $k+j<2(j-1)+1=2j-1$, czyli $j>k+1$.
Dla $X=S^0$ otrzymujemy $\Sigma^nS^0\cong S^n$
i ten sam rachunek daje $n\ge k+2$.
\end{proof}

\begin{wniosek}[Zerowy stabilny pień sfer]
$\pi_0^{\mathrm s}\cong\Z$.
\end{wniosek}
\begin{proof}
Dla $n\ge2$ mamy $\pi_n(S^n)\cong\Z$,
gdzie izomorfizm wyznacza stopień odwzorowania.
Zawieszenie zachowuje stopień, co można sprawdzić
na homologii fundamentalnej i izomorfizmie zawieszenia
$\widetilde H_n(S^n)\cong\widetilde H_{n+1}(S^{n+1})$.
Dlatego system stabilizujący jest układem identyczności
na $\Z$.
\end{proof}

\section{Niezerowy stabilny element Hopfa}
\begin{lemat}[Rozszczepienie zawieszonego iloczynu]
\label{lem:sigma-product-v8}
Dla punktowanych CW-kompleksów $X,Y$ istnieje naturalna
równoważność homotopijna
$\Sigma(X\times Y)\simeq
\Sigma X\vee\Sigma Y\vee\Sigma(X\wedge Y)$.
W szczególności zawieszenie odwzorowania dołączającego
$S^{p+q-1}\to S^p\vee S^q$ górnej komórki
$S^p\times S^q$, dla $p,q\ge2$, jest nulhomotopijne.
\end{lemat}
\begin{proof}
Ponieważ włączenia punktów bazowych są kofibracjami,
$X\vee Y\hookrightarrow X\times Y$ jest kofibracją,
a jej przestrzeń ilorazowa jest homeomorficzna z
$X\wedge Y$.
Po zawieszeniu włączenie
$\Sigma(X\vee Y)\to\Sigma(X\times Y)$
ma lewy homotopijny odwrotny: tworzą go zawieszone
projekcje na $X$ i $Y$ połączone operacją uszczypnięcia
na zawieszeniu. To rozszczepia zawieszoną kofibrację
w kategorii homotopijnej. W otrzymanym rozkładzie
trzeci składnik ma typ $\Sigma(X\wedge Y)$,
gdyż iloraz pary ma ten typ i zawieszone włączenie
jest rozszczepione. Dla dwóch sfer część pochodząca
z górnej komórki jest więc oddzielnym składnikiem
$S^{p+q+1}$, a zawieszone odwzorowanie dołączające jest
nulhomotopijne.
\end{proof}

\begin{twierdzenie}[Pierwszy stabilny pień]
\label{thm:stable-eta-v8}
Niech $\eta:S^3\to S^2$ będzie odwzorowaniem Hopfa o
$H(\eta)=1$. Wówczas zawieszenie $E\eta$
reprezentuje niezerowy element rzędu $2$ w
$\pi_4(S^3)$, a
$\pi_1^{\mathrm s}\cong\Z/2\Z$.
\end{twierdzenie}
\begin{proof}
Wiemy z poprzednich rozdziałów, że
$\pi_3(S^2)\cong\Z$ z generatorem $\eta$.
Z twierdzenia Freudenthala otrzymujemy surjekcję
$E:\pi_3(S^2)\twoheadrightarrow\pi_4(S^3)$.
Z Twierdzenia~\ref{thm:whitehead-square-v7} mamy
$[\iota_2,\iota_2]=\pm2\eta$.
Lemat~\ref{lem:sigma-product-v8} mówi, że zawieszenie
iloczynu Whiteheada jest zerowe, więc $2E\eta=0$.
Wynika stąd, że $\pi_4(S^3)$ jest albo zerowa,
albo izomorficzna z $\Z/2\Z$.

Załóżmy, że $E\eta=0$. Wtedy stożek odwzorowania
$\Sigma\eta:S^4\to S^3$ ma typ
$S^3\vee S^5$. Z drugiej strony zawieszenie zachowuje
konstrukcję stożka odwzorowania, zatem
$C_{\Sigma\eta}\simeq\Sigma C_\eta
\simeq\Sigma\C P^2$.
Twierdzenie~\ref{thm:sigma-cp2-v8} udowodniło,
że te przestrzenie nie mogą być równoważne:
$\operatorname{Sq}^2$ jest nietrywialne na pierwszej
i zerowe na drugiej. Sprzeczność dowodzi $E\eta\ne0$.
Stąd $\pi_4(S^3)\cong\Z/2\Z$.
Dla $n\ge3$ grupy $\pi_{n+1}(S^n)$ są już w zakresie
stabilnym z Twierdzenia~\ref{thm:stable-range-v8},
więc wyznaczają $\pi_1^{\mathrm s}\cong\Z/2\Z$.
\end{proof}

\begin{przyklad}[Jawny układ stabilizacji pierwszego pnia]
\label{prz:eta-stabilizacja-v8}
Kolejne zawieszenia odwzorowania Hopfa dają układ
\[
 \underbrace{\pi_3(S^2)}_{\Z}
 \xrightarrow{\ E\ }
 \underbrace{\pi_4(S^3)}_{\Z/2}
 \xrightarrow{\ E\ }
 \underbrace{\pi_5(S^4)}_{\Z/2}
 \xrightarrow{\ E\ }\cdots.
\]
Pierwsza strzałka jest redukcją modulo $2$
przy identyfikacji generatorów przez $\eta$:
z dowodu Twierdzenia~\ref{thm:stable-eta-v8}
wiemy, że $E\eta$ ma dokładnie rząd $2$
i generuje grupę po prawej, więc
$\ker E=2\Z$.
Pozostałe strzałki od poziomu $S^3$ są
izomorfizmami na mocy Freudenthala:
dla $S^n$ w stopniu $n+1$ mamy
$r=n-1$, $q=n+1$, a nierówność
$q<2r+1$ sprowadza się do $n>2$.
W granicy prostej otrzymujemy
$\pi_1^{\mathrm s}=\Z/2$.

Dla porównania układ stopni
$\pi_n(S^n)\cong\Z$ ma wszystkie
strzałki zawieszenia równe identyczności
po uzgodnieniu orientacji; jego granicą
jest $\pi_0^{\mathrm s}\cong\Z$.
Tak więc \emph{pierwsze dwa} stabilne
pnie sfer mają różną naturę:
pierwszy wykrywa stopień całkowity,
a drugi niezerowe zawieszenie odwzorowania Hopfa.
\end{przyklad}

\begin{figure}[!htbp]
\centering
\begin{tikzpicture}[>=Stealth,font=\small]
 \node[draw=manifoldblue!75!black,rounded corners=4pt,
       fill=manifoldblue!8,minimum width=2.5cm,
       minimum height=1.1cm,align=center] (a) at (0,0)
       {$\pi_3(S^2)$\\$\Z$};
 \node[draw=manifoldgreen!75!black,rounded corners=4pt,
       fill=manifoldgreen!8,minimum width=2.5cm,
       minimum height=1.1cm,align=center] (b) at (4.25,0)
       {$\pi_4(S^3)$\\$\Z/2$};
 \node[draw=manifoldgreen!75!black,rounded corners=4pt,
       fill=manifoldgreen!8,minimum width=2.5cm,
       minimum height=1.1cm,align=center] (c) at (8.5,0)
       {$\pi_5(S^4)$\\$\Z/2$};
 \draw[->,manifoldred,thick] (a)--(b)
       node[midway,above,font=\scriptsize] {$E:\ \bmod 2$};
 \draw[->,manifoldblue!75!black,thick] (b)--(c)
       node[midway,above,font=\scriptsize] {$E:\ \cong$};
 \node[below=7pt,align=center] at (4.25,-.55)
       {Od $S^3$ wszystkie kolejne\\
        zawieszenia są izomorfizmami};
\end{tikzpicture}
\caption{Stabilizacja pierwszego dodatniego pnia sfer:
pierwsza strzałka redukuje liczbę Hopfa modulo $2$,
a następne są izomorfizmami.}
\label{fig:eta-stabilizacja-v8}
\end{figure}

\begin{uwaga}[Co wykazał ten dowód?]
Wykorzystaliśmy tylko: homotopię fibracji Hopfa,
równanie Whiteheada, rozszczepienie zawieszonego
iloczynu, Freudenthala oraz stabilność kwadratów
Steenroda. Nie potrzebowaliśmy tabel wszystkich
niestabilnych grup homotopii sfer. Najtrudniejszy
zewnętrzny składnik argumentu to stabilność operacji
Steenroda, jawnie wskazana w poprzednim rozdziale.
\end{uwaga}

\section{Stabilne odwzorowania między skończonymi kompleksami}
\begin{definicja}[Klasa stabilna odwzorowania]
Dla skończonych punktowanych CW-kompleksów $X,Y$
przyjmujemy
\[
 [X,Y]_{\mathrm s}:=\varinjlim_{j\ge0}
 [\Sigma^jX,\Sigma^jY]_*,
\]
gdzie $[-,-]_*$ oznacza punktowane klasy homotopii.
Nie zakładamy, że jest to już grupa dla $j=0$;
po co najmniej jednym zawieszeniu klasy tworzą
grupy dzięki strukturze ko-H przestrzeni źródłowej.
Po dalszym zawieszeniu otrzymujemy grupy abelowe,
co można uzasadnić dwiema zgodnymi operacjami
składania po dwóch współrzędnych zawieszenia
i argumentem Eckmanna--Hiltona.
Otrzymana granica prosta ma więc naturalną
strukturę grupy abelowej.
\end{definicja}

\begin{stwierdzenie}[Stabilizacja klas odwzorowań]
Niech $X,Y$ będą skończonymi punktowanymi
CW-kompleksami, a $Y$ niech będzie łukowo spójna.
Dla dostatecznie dużych $j$ zawieszenie
$[\Sigma^jX,\Sigma^jY]_*
\to[\Sigma^{j+1}X,\Sigma^{j+1}Y]_*$
jest bijekcją.
\end{stwierdzenie}
\begin{proof}
Niech $d=\dim X$. Dla $j\ge1$ przestrzeń
$Z_j=\Sigma^jY$ jest co najmniej $j$-spójna,
a $W_j=\Sigma^jX$ ma wymiar nie większy niż $d+j$.
Przez adjunkcję zawieszenia i pętli rozpatrywane odwzorowanie
identyfikuje się z odwzorowaniem
$[W_j,Z_j]_*\to[W_j,\Omega\Sigma Z_j]_*$
indukowanym przez jednostkę $u_j:Z_j\to\Omega\Sigma Z_j$.
Twierdzenie Freudenthala zapewnia, że $u_j$
indukuje izomorfizmy grup homotopii w stopniach
$q<2j+1$ oraz surjekcję w stopniu $q=2j+1$.
Jest zatem $(2j+1)$-spójna.

Używamy tu następującej komórkowej postaci zasady
przedłużania homotopii: jeśli $u:Z\to Z'$ jest
$N$-spójnym odwzorowaniem przestrzeni punktowanych,
to znaczy indukuje izomorfizmy $\pi_i$ dla
$i<N$ oraz surjekcję dla $i=N$, to
$[W,Z]_*\to[W,Z']_*$ jest bijekcją dla
punktowanego CW-kompleksu $W$ o $\dim W<N$.
Istnienie podniesienia do homotopii wynika ze
znikania odpowiednich względnych grup homotopii
na kolejnych komórkach $W$; dla różnowartościowości
stosuje się ten sam argument do homotopii na
$W\times I$ z ustalonymi końcami (w wymiarze
co najwyżej $\dim W+1\le N$).
Dla $j\ge d+1$ mamy $d+j<2j+1$.
Stąd zawieszenie jest bijekcją w każdym takim kroku,
co dowodzi stabilizacji.
\end{proof}
\begin{uwaga}
Istotnym składnikiem jest \emph{komórkowa zasada
przedłużania} dla $N$-spójnego odwzorowania, a nie sam fakt,
że indukuje ona izomorfizmy na pewnych grupach
homotopii. Przedstawiony argument przez szkielety
wskazuje dokładnie, gdzie stosuje się ten wynik.
\end{uwaga}

\section{Widma: dlaczego stabilizacji nie opisuje pojedyncza przestrzeń?}
\begin{definicja}[Widmo ciągowe]
\emph{Widmem ciągowym} nazywamy rodzinę
punktowanych przestrzeni $E_0,E_1,\ldots$ i
odwzorowań struktury $\sigma_n:\Sigma E_n\to E_{n+1}$.
Jeżeli dołączone odwzorowania $E_n\to\Omega E_{n+1}$
są słabymi równoważnościami dla wszystkich $n$,
mówimy o \emph{$\Omega$-widmie}.
Dla $k\in\Z$ stabilne grupy homotopii widma
definiujemy przez
$\pi_k(E)=\varinjlim_{n\ge\max(0,1-k)}\pi_{n+k}(E_n)$,
z odwzorowaniami indukowanymi przez zawieszenie i struktury.
\end{definicja}

\begin{przyklad}[Widmo sferyczne]
Widmo $\mathbb S$ ma $\mathbb S_n=S^n$ i odwzorowania
struktury $\Sigma S^n\cong S^{n+1}$.
Jego grupy homotopii to stabilne pnie sfer:
$\pi_k(\mathbb S)=\pi_k^{\mathrm s}$ dla $k\ge0$.
To widmo nie jest $\Omega$-widmem, gdyż
$S^n\not\simeq\Omega S^{n+1}$ w ogólności.
\end{przyklad}

\begin{przyklad}[Widmo Eilenberga--Mac Lane'a]
Niech $A$ będzie grupą abelową. Wybierzmy
$HA_n=K(A,n)$ dla $n\ge0$, gdzie $K(A,0)$
traktujemy jako przestrzeń dyskretną z wyróżnionym zerem.
Standardowa słaba równoważność
$K(A,n)\simeq\Omega K(A,n+1)$ wyznacza odwzorowania
struktury $\Sigma K(A,n)\to K(A,n+1)$.
Dla $n\ge1$ jest to konsekwencja obliczenia
$\pi_j(\Omega K(A,n+1))\cong\pi_{j+1}(K(A,n+1))$.
W stopniu $n=0$ przestrzeń $\Omega K(A,1)$
ma składowe łukowe indeksowane przez
$\pi_1(K(A,1))\cong A$, a każda składowa
ma zerowe dodatniostopniowe grupy homotopii;
jest więc słabo równoważna dyskretnej przestrzeni $A$.
Wtedy $HA$ jest $\Omega$-widmem w sensie słabej
równoważności, a jego grupy stabilne wynoszą
$\pi_k(HA)=A$ dla $k=0$ i $0$ dla $k\ne0$.
Dla $k=0$ na każdym poziomie $n\ge1$ mamy
$\pi_n(K(A,n))\cong A$. Ważne jest również
sprawdzenie strzałek układu: są to złożenia
homomorfizmu zawieszenia i homomorfizmu
indukowanego przez odwzorowanie struktury.
Po zastosowaniu adjunkcji odpowiadają one odwzorowaniom
$\pi_n(K(A,n))\to\pi_n(\Omega K(A,n+1))$,
które są izomorfizmami, ponieważ wybrane
odwzorowania sprzężone są słabymi równoważnościami.
Zatem granicą jest rzeczywiście $A$,
a nie tylko grupa o nieznanym działaniu strzałek.
Dla $k\ne0$ wszystkie późniejsze grupy
$\pi_{n+k}(K(A,n))$ są zerowe.
\end{przyklad}

\begin{twierdzenie}[Kohomologia jako stabilna teoria]
\label{thm:HA-rep-v8}
Dla punktowanego CW-kompleksu $X$, grupy abelowej $A$
i liczby całkowitej $m\ge1$ zachodzi naturalna bijekcja
$\widetilde H^m(X;A)\cong[X,K(A,m)]_*$.
Po zawieszeniu otrzymujemy zgodne identyfikacje
$\widetilde H^{m+j}(\Sigma^jX;A)
\cong[\Sigma^jX,K(A,m+j)]_*$.
Wszystkie te identyfikacje są zgodne z izomorfizmem
zawieszenia w zredukowanej kohomologii.
\end{twierdzenie}
\begin{proof}
Pierwsza bijekcja jest twierdzeniem o reprezentowalności
kohomologii przez $K(A,m)$ z wcześniejszego rozdziału.
Zredukowany izomorfizm zawieszenia wynika z długiego
ciągu dokładnego pary $(CX,X)$ i ściągalności stożka.
Po stronie przestrzeni Eilenberga--Mac Lane'a
przechodzi on przez adjunkcję $\Sigma\dashv\Omega$ i
słabą równoważność
$K(A,m)\simeq\Omega K(A,m+1)$.
Naturalność wynika z naturalności tych dwóch konstrukcji.
\end{proof}

\begin{uwaga}[Od kohomologii do teorii reprezentowanych]
Pełne stwierdzenie spektralne ma postać
$\widetilde H^m(X;A)\cong[\Sigma^\infty X,\Sigma^mHA]$
w stabilnej kategorii homotopii. Tutaj
$\Sigma^\infty X$ oznacza \emph{widmo zawieszeń}
(o $n$-tym poziomie $\Sigma^nX$), a nawias
$[E,F]$ zbiór stabilnych klas odwzorowań widm. Jego dowód wymaga
konstrukcji odwzorowań widm, przejścia do modeli
homotopijnych oraz twierdzenia o reprezentowalności.
W obecnym skrypcie podajemy dokładną wersję dla
$K(A,m)$, a zapis spektralny wskazuje kierunek
dalszego uogólnienia.
\end{uwaga}

\section{Czego nie można odczytać z samej tabeli grup?}
Znajomość wszystkich grup homologii $H_*(X)$ nie
wyznacza stabilnego typu homotopijnego: zawieszona
$\C P^2$ i bukiet sfer mają identyczne homologie,
ale różnią się działaniem $\operatorname{Sq}^2$.
Podobnie same grupy homotopii mogą nie wyznaczać
niestabilnego typu homotopijnego: potrzebne są
niezmienniki Postnikowa.

Kolejnymi krokami są ciąg spektralny Adamsa,
wykorzystujący algebrę Steenroda i funktory $\operatorname{Ext}$,
oraz twierdzenie Pontriagina--Thoma wiążące
stabilną homotopię z kobordyzmami. Dopiero po ich
przygotowaniu sensowne jest systematyczne
obliczanie dalszych stabilnych pni sfer.

